\documentclass[a4j, 9pt]{ltjsarticle}

\usepackage{mypackage}
\usepackage{subfiles}

\begin{document}

% 目次の出力
\tableofcontents
\clearpage

\part{基礎編}

  \section{数学記号一覧}

  \newpage

  \subfile{sections/logic.tex}

  \newpage

  \subfile{sections/set.tex}

  \newpage

  \subfile{sections/mapping.tex}

  \newpage
  
  \subfile{sections/symmetry.tex}

  \newpage

  \subfile{sections/combinatorics.tex}

  \newpage
  
\part{発展編}
  基礎編では数学の基礎、使いこなせると役に立つ道具類を紹介してきた。
  本パートでは、少々込み入った専門的な知識を確認していくことにしよう。
  特に代数学では、代表的な\textbf{群論}、\textbf{環論}、\textbf{体}を主として述べることにする。

  \section{代数学}
    

\end{document}
